% Reconstruccion editable del pensum de Ingenieria Biomedica.
% Todo el contenido es texto y tablas LaTeX; no se incrusta el PDF original.
\documentclass[10pt,letterpaper]{article}

\usepackage[spanish,es-nodecimaldot]{babel}
\usepackage[utf8]{inputenc}
\usepackage[T1]{fontenc}
\usepackage{lmodern}
\usepackage[margin=1.6cm]{geometry}
\usepackage[table]{xcolor}
\usepackage{array,longtable,booktabs,tabularx}
\usepackage{enumitem}
\usepackage{microtype}
\usepackage[hidelinks]{hyperref}

\definecolor{unitecblue}{HTML}{005A9C}
\definecolor{lightblue}{HTML}{EAF3F8}
\definecolor{bmdgreen}{HTML}{168447}
\setlength{\parindent}{0pt}
\setlength{\parskip}{5pt}
\setlist[itemize]{leftmargin=1.3em,itemsep=2pt,topsep=3pt,label=\textcolor{unitecblue}{\guillemotright}}
\renewcommand{\arraystretch}{1.15}
\newcommand{\uv}[1]{#1 U.V.}
\newcommand{\bmdcourse}[2]{\textcolor{bmdgreen}{\bfseries #1} & \textcolor{bmdgreen}{#2}}
\newcommand{\coursecode}[1]{\texttt{#1}}

\begin{document}

{\color{unitecblue}\Huge\bfseries Ingeniería Biomédica}\hfill
{\Large Facultad de Ingeniería}

\vspace{4pt}\hrule\vspace{10pt}

\section*{Perfil profesional}

El egresado de Ingeniería Biomédica es un profesional de referencia nacional e
internacional, con conocimientos, habilidades y competencias que le permiten
desempeñarse en los diversos campos de aplicación de la carrera. Con un enfoque
en la gestión de la tecnología médica, tiene una clara orientación hacia las
necesidades del sector salud y de la sociedad, planteando soluciones tecnológicas
con innovación e investigación. Aplica métodos de ingeniería y tecnología en el
diseño y evaluación de áreas hospitalarias bajo el cumplimiento de normativas
vigentes; desarrolla e implementa dispositivos y tecnologías médicas con el fin
de brindar alternativas viables para mejorar la prestación de servicios de salud,
con sentido ético y compromiso social.

\section*{Competencias específicas}

\begin{itemize}
  \item Utilizar métodos de ingeniería y tecnología para solventar necesidades en la prestación de servicios de salud.
  \item Aplicar fundamentos matemáticos, físicos y científicos para identificar, formular y resolver problemas de ingeniería biomédica y clínica.
  \item Interpretar y utilizar los fundamentos de los sistemas fisiológicos humanos, tanto a nivel estructural como funcional, y las patologías más relevantes para aplicar soluciones de ingeniería biomédica que satisfagan la necesidad clínica y social.
  \item Gestionar y analizar la instrumentación biomédica para diagnóstico, monitorización y terapia que interactúa con el paciente, mediante la comprensión de principios matemáticos, físicos, mecánicos, electrónicos y de ingeniería.
  \item Evaluar la infraestructura y tecnología hospitalaria requerida desde el punto de vista de su efectividad y valoración para uso médico.
  \item Identificar, operar y diseñar sistemas de información utilizados en medicina.
  \item Desarrollar e implementar algoritmos para resolver problemas de ingeniería biomédica, considerando la normativa vigente y las responsabilidades éticas y sociales.
  \item Definir y utilizar sensores, actuadores y sistemas de adquisición de señales biomédicas para evaluar o diseñar dispositivos médicos.
  \item Aplicar metodologías de mejora o desarrollo de dispositivos médicos bajo experimentación apropiada.
  \item Diseñar y llevar a cabo proyectos tecnológicos en el ámbito de la ingeniería biomédica.
  \item Adquirir nuevos conocimientos de ingeniería mediante métodos efectivos de aprendizaje y aplicarlos para resolver necesidades médicas.
  \item Dirigir, comunicarse e interactuar efectivamente con audiencias afines a la Ingeniería Biomédica y las ciencias de la salud.
\end{itemize}

\section*{Campo laboral}

\begin{itemize}
  \item Departamentos de ingeniería biomédica de hospitales y clínicas.
  \item Empresas distribuidoras o representantes de fabricantes de tecnología médica.
  \item Empresas prestadoras de servicio o soporte técnico de tecnología o infraestructura médica.
  \item Instituciones gubernamentales de salud.
  \item Consultoría biomédica.
  \item Fundaciones y brigadas médicas.
  \item Docencia e investigación.
  \item Empresa propia de servicios biomédicos.
\end{itemize}

\clearpage
\section*{Plan de estudios 2020}

\small
\rowcolors{2}{lightblue}{white}
\begin{longtable}{>{\bfseries}p{1.7cm}p{8.4cm}>{\centering\arraybackslash}p{1.3cm}>{\centering\arraybackslash}p{2.0cm}}
\rowcolor{unitecblue}\color{white}Código & \color{white}Asignatura & \color{white}U.V. & \color{white}Teoría--práctica \\
\endfirsthead
\rowcolor{unitecblue}\color{white}Código & \color{white}Asignatura & \color{white}U.V. & \color{white}Teoría--práctica \\
\endhead
MAT101 & Introducción al Álgebra & 4 & 4--0 \\
ESP103 & Comunicación Oral y Escrita & 2 & 2--0 \\
\bmdcourse{BMD101}{Introducción a la Ingeniería Biomédica} & 2 & 1--1 \\
LCC104 & Ofimática Avanzada & 4 & 4--0 \\
MAT102 & Álgebra & 4 & 4--0 \\
HIS101 & Historia de Honduras & 3 & 3--0 \\
SOC101 & Sociología & 3 & 3--0 \\
EIE1 & Electiva de Idioma Extranjero I & 3 & 3--0 \\
MAT103 & Geometría y Trigonometría & 4 & 4--0 \\
ART/DEP & Electiva de Arte o Deporte & 3 & 3--0 \\
FIL101 & Filosofía & 3 & 3--0 \\
EIE2 & Electiva de Idioma Extranjero II & 3 & 3--0 \\
MAT109 & Cálculo I & 4 & 4--0 \\
MAT105 & Álgebra Lineal & 4 & 4--0 \\
QUI101 & Química General & 4 & 4--0 \\
EIE3 & Electiva de Idioma Extranjero III & 3 & 3--0 \\
MAT210 & Cálculo II & 4 & 4--0 \\
MEC301 & Dibujo Técnico & 3 & 2--1 \\
TLL310 & Ética y Ciudadanía & 2 & 1--1 \\
EIE4 & Electiva de Idioma Extranjero IV & 3 & 3--0 \\
MAT203 & Ecuaciones Diferenciales & 4 & 4--0 \\
FIS201 & Física I & 4 & 4--0 \\
MAT301 & Estadística Matemática I & 4 & 4--0 \\
TLL314 & Metodología de Investigación & 1 & 0--1 \\
MAT209 & Variable Compleja & 3 & 3--0 \\
FIS202 & Física II & 4 & 4--0 \\
\bmdcourse{BMD303}{Anatomía y Fisiología I} & 4 & 4--0 \\
ADM102 & Administración I & 4 & 4--0 \\
SEL311 & Análisis de Circuitos Eléctricos en D.C. & 3 & 3--0 \\
EMP401 & Generación de Empresas I & 3 & 3--0 \\
FIS203 & Física III & 4 & 4--0 \\
\bmdcourse{BMD305}{Anatomía y Fisiología II} & 4 & 3--1 \\
SEL312 & Análisis de Circuitos Eléctricos en A.C. & 3 & 3--0 \\
EMP402 & Generación de Empresas II & 3 & 3--0 \\
\bmdcourse{BMD301}{Termodinámica} & 3 & 3--0 \\
\bmdcourse{BMD304}{Instalaciones Hospitalarias I} & 4 & 3--1 \\
SEL313 & Análisis de Circuitos Electrónicos & 4 & 3--1 \\
\bmdcourse{BMD308}{Seminarios Clínicos} & 3 & 3--0 \\
\bmdcourse{BMD306}{Instalaciones Hospitalarias II} & 3 & 3--0 \\
\bmdcourse{BMD203}{Regulación Sanitaria} & 3 & 3--0 \\
SEL314 & Amplificadores Electrónicos y A.O. & 4 & 4--0 \\
SEL308 & Ingeniería de Control & 3 & 3--0 \\
\bmdcourse{BMD402}{Instalaciones Hospitalarias III} & 3 & 3--0 \\
\bmdcourse{BMD307}{Ingeniería Clínica} & 4 & 3--1 \\
CCC214 & Programación C++ & 4 & 4--0 \\
MEC306 & Sistemas de Gases Medicinales & 3 & 3--0 \\
\bmdcourse{BMD311}{Bioinstrumentación y Tecnologías Biomédicas I} & 4 & 4--0 \\
\bmdcourse{BMD309}{Sensores y Actuadores} & 3 & 3--0 \\
\bmdcourse{BMD416}{Procesamiento de Señales Biomédicas} & 3 & 3--0 \\
\bmdcourse{BMD417}{Tecnologías de Laboratorio Clínico} & 4 & 3--1 \\
\bmdcourse{BMD310}{Bioinstrumentación y Tecnologías Biomédicas II} & 3 & 3--0 \\
SEL405 & Teoría del Control Digital & 3 & 3--0 \\
\bmdcourse{BMD418}{Sistemas de Imágenes Médicas} & 3 & 3--0 \\
\bmdcourse{BMD312}{Biomecánica y Desarrollo de Biodispositivos} & 4 & 4--0 \\
\bmdcourse{BMD419}{Análisis de Dispositivos Médicos} & 4 & 4--0 \\
EFEBMD1 & Electiva de Formación Específica I & 3 & 2--1 \\
\bmdcourse{BMD420}{Informática Médica} & 3 & 2--1 \\
\bmdcourse{BMD421}{Procesamiento de Imágenes Médicas} & 4 & 4--0 \\
IND432 & Sistemas de Gestión de la Innovación y Tecnología & 4 & 4--0 \\
EFEBMD2 & Electiva de Formación Específica II & 3 & 3--0 \\
\bmdcourse{BMD428}{Proyecto de Investigación\textsuperscript{**}} & 4 & 0--4 \\
EFEBMD3 & Electiva de Formación Específica III & 4 & 4--0 \\
\bmdcourse{BMD503 / BMD504}{Elegir una: Proyecto de Graduación o Práctica Profesional\textsuperscript{***}} & 4 & 0--4 \\
\bottomrule
\end{longtable}

\normalsize
\section*{Relaciones de prerrequisito}

La flecha significa que la asignatura de la izquierda debe aprobarse antes de
cursar la asignatura de la derecha. Las cadenas agrupan las rutas de precedencia
del flujograma para que puedan leerse y editarse como texto.

\small
\begin{longtable}{>{\raggedleft\arraybackslash}p{4.1cm}>{\centering\arraybackslash}p{1cm}p{8.2cm}}
\toprule
Requisito & & Asignatura habilitada \\
\midrule
\endfirsthead
\toprule
Requisito & & Asignatura habilitada \\
\midrule
\endhead
\coursecode{MAT101} & $\longrightarrow$ & \coursecode{MAT102} \\
\coursecode{MAT102} & $\longrightarrow$ & \coursecode{MAT103}, \coursecode{MAT105} \\
\coursecode{MAT103} & $\longrightarrow$ & \coursecode{MAT109} \\
\coursecode{MAT109} & $\longrightarrow$ & \coursecode{MAT210}, \coursecode{FIS201} \\
\coursecode{MAT210} & $\longrightarrow$ & \coursecode{MAT203} \\
\coursecode{FIS201} & $\longrightarrow$ & \coursecode{FIS202} \\
\coursecode{FIS202} & $\longrightarrow$ & \coursecode{FIS203}, \coursecode{SEL311} \\
\coursecode{FIS203} & $\longrightarrow$ & \textcolor{bmdgreen}{\coursecode{BMD301}} \\
\textcolor{bmdgreen}{\coursecode{BMD101}} & $\longrightarrow$ & \textcolor{bmdgreen}{\coursecode{BMD303}} \\
\textcolor{bmdgreen}{\coursecode{BMD303}} & $\longrightarrow$ & \textcolor{bmdgreen}{\coursecode{BMD305}} \\
\textcolor{bmdgreen}{\coursecode{BMD305}} & $\longrightarrow$ & \textcolor{bmdgreen}{\coursecode{BMD308}} \\
\coursecode{SEL311} & $\longrightarrow$ & \coursecode{SEL312} \\
\coursecode{SEL312} & $\longrightarrow$ & \coursecode{SEL313} \\
\coursecode{SEL313} & $\longrightarrow$ & \coursecode{SEL314} \\
\coursecode{SEL314} & $\longrightarrow$ & \textcolor{bmdgreen}{\coursecode{BMD311}}, \coursecode{SEL405}, \coursecode{SEL407} \\
\textcolor{bmdgreen}{\coursecode{BMD304}} & $\longrightarrow$ & \textcolor{bmdgreen}{\coursecode{BMD306}} \\
\textcolor{bmdgreen}{\coursecode{BMD306}} & $\longrightarrow$ & \textcolor{bmdgreen}{\coursecode{BMD402}} \\
\textcolor{bmdgreen}{\coursecode{BMD402}} & $\longrightarrow$ & \textcolor{bmdgreen}{\coursecode{BMD307}} \\
\textcolor{bmdgreen}{\coursecode{BMD311}} & $\longrightarrow$ & \textcolor{bmdgreen}{\coursecode{BMD310}} \\
\textcolor{bmdgreen}{\coursecode{BMD309}} & $\longrightarrow$ & \textcolor{bmdgreen}{\coursecode{BMD310}} \\
\coursecode{CCC214} & $\longrightarrow$ & \textcolor{bmdgreen}{\coursecode{BMD416}} \\
\textcolor{bmdgreen}{\coursecode{BMD418}} & $\longrightarrow$ & \textcolor{bmdgreen}{\coursecode{BMD421}} \\
\textcolor{bmdgreen}{\coursecode{BMD312}} & $\longrightarrow$ & \textcolor{bmdgreen}{\coursecode{BMD419}} \\
\coursecode{TLL314} & $\longrightarrow$ & \textcolor{bmdgreen}{\coursecode{BMD428}} \\
\textcolor{bmdgreen}{\coursecode{BMD428}} & $\longrightarrow$ & \textcolor{bmdgreen}{\coursecode{BMD503 / BMD504}} \\
\bottomrule
\end{longtable}

\normalsize
\section*{Electivas de formación específica}

\small
\begin{tabularx}{\textwidth}{>{\bfseries}p{1.7cm}X>{\centering\arraybackslash}p{1.4cm}p{2.1cm}}
\toprule
Código & Asignatura & U.V. & Requisito \\
\midrule
SEL407 & Electrónica de Potencia & 4 & SEL314 \\
\textcolor{bmdgreen}{\bfseries BMD422} & \textcolor{bmdgreen}{Tecnologías de Odontología} & 4 & \textcolor{bmdgreen}{BMD417} \\
SEL402 & Microprocesadores I & 4 & SEL405 \\
\textcolor{bmdgreen}{\bfseries BMD423} & \textcolor{bmdgreen}{Gestión de Tecnología Médica} & 4 & IND432 \\
\textcolor{bmdgreen}{\bfseries BMD424} & \textcolor{bmdgreen}{Tecnologías de Oftalmología} & 4 & \textcolor{bmdgreen}{BMD417} \\
\textcolor{bmdgreen}{\bfseries BMD425} & \textcolor{bmdgreen}{Electrónica y Fotónica en Dispositivos Médicos} & 4 & SEL405 \\
\textcolor{bmdgreen}{\bfseries BMD426} & \textcolor{bmdgreen}{Infraestructura Hospitalaria} & 4 & IND432 \\
\textcolor{bmdgreen}{\bfseries BMD427} & \textcolor{bmdgreen}{Tecnologías de Laboratorios de Patología} & 4 & \textcolor{bmdgreen}{BMD417} \\
IND434 & Normalización y Auditoría de la Calidad & 4 & IND432 \\
\bottomrule
\end{tabularx}

\normalsize
\section*{Requisitos especiales}

\begin{itemize}
  \item \textsuperscript{*} Requisito: 172 U.V. aprobadas.
  \item \textsuperscript{**} Requisito: 202 U.V. aprobadas e índice académico mayor o igual a 70\%.
  \item \textsuperscript{***} Requisito: 210 U.V. aprobadas, índice académico mayor o igual a 70\% y cumplimiento del requisito de graduación de competencias idiomáticas.
\end{itemize}

\section*{Distribución por bloque de conocimiento}

\begin{tabularx}{\textwidth}{Xrrr}
\toprule
Bloque del conocimiento & Asignaturas & U.V. & \% \\
\midrule
Formación General y Complementaria & 15 & 45 & 21 \\
Ciencias Básicas & 16 & 61 & 29 \\
Tecnologías Básicas & 8 & 31 & 14 \\
Ingeniería Biomédica Aplicada & 22 & 69 & 32 \\
Formación Integradora & 2 & 8 & 4 \\
\midrule
\textbf{Total} & \textbf{63} & \textbf{214} & \textbf{100} \\
\bottomrule
\end{tabularx}

\vfill
\footnotesize Fuente reconstruida del documento institucional ``Flujograma Ing. Biomédica 2020'', versión VM09042025V.\\
\url{https://www.unitec.edu}

\end{document}
